% ============================================================================
% 6. DISCUSSION — Journal version (expanded)
% ============================================================================
\section{Discussion}
\label{sec:discussion}

\subsection{The Retrieval--Reasoning Decomposition}

Our results reveal a clean decomposition of graph reasoning performance
into two independent failure modes.

\textbf{Retrieval-limited subtypes} (\texttt{hop\_count},
\texttt{direct\_fk}, \texttt{isolation}, \texttt{connected}) achieve
$\geq$93\% EM with Tool-Use because the answer requires only one or
two graph lookups.  Static baselines fail on these subtypes
\emph{not because the reasoning is complex}, but because the relevant
graph fragment is missing from the injected context.  This is confirmed
by the Oracle's near-perfect scores: once the correct information is
available, both models can interpret it accurately.

\textbf{Reasoning-limited subtypes} (\texttt{combined\_impact},
\texttt{membership}, \texttt{transitive}) remain hard even with tools.
Even with Tool-Use providing graph algorithm access,
\texttt{combined\_impact} achieves only 12\% EM ($n=69$).  The LLM
struggles to compose multiple tool outputs across phases, confirming
a ceiling for prompt-based graph reasoning.  Notably, RC
(5~code execution rounds) shows a similar ceiling to TU (3~tool calls),
suggesting the bottleneck is compositional reasoning itself rather
than interaction budget.

\subsection{The Tool Necessity-Insufficiency Paradox}

Tier~2 reveals a striking paradox: tools are simultaneously
\emph{essential} and \emph{insufficient}.

On data-dependent subtypes (\texttt{cascade\_count}: 100\% vs.\ 0\%;
\texttt{row\_provenance}: 92\% vs.\ 0\%), tools are the difference
between perfect performance and complete failure.  These subtypes
require traversing the 18MB row-level lineage map---far exceeding any
context window.

Yet on reasoning-intensive subtypes (\texttt{multi\_hop\_trace}: 27\%;
\texttt{shared\_source}: 23\%), tools provide \emph{no} advantage over
static baselines.  These subtypes require multi-hop graph traversal
and set-theoretic reasoning (finding common ancestors of two mart
tables)---capabilities that tools provide but LLMs cannot effectively
orchestrate.

\subsection{Implications for Hybrid Architectures}

The complementarity analysis confirms that Tool-Use and Graph-Augmented
fail on largely the \emph{same} questions: their union improves overall
EM by only 6 percentage points, leaving 35\% of questions unsolved by
either approach.  This suggest that the bottleneck is not retrieval
\emph{or} tool access, but the LLM's ability to perform compositional
structural reasoning.

This points toward hybrid architectures combining:
\begin{enumerate}
  \item \textbf{Tool-augmented retrieval} for data-access subtypes
    where simple lookups suffice.
  \item \textbf{Learned structural representations} (e.g., GNN encoders)
    that can capture multi-hop topology patterns beyond what prompt-based
    reasoning can achieve.
  \item \textbf{Compositional reasoning modules} that decompose complex
    graph queries into sequences of simpler operations.
\end{enumerate}

\subsection{Obfuscation and Memorization}

The obfuscation results provide a clean diagnostic for memorization.
Static baselines lose 15--32\% EM under obfuscation, confirming partial
reliance on memorized schema name associations.  Tool-Use's near-invariance
($\Delta{=}{-}3.4$\% Gemini, ${-}4.2$\% DeepSeek) demonstrates that
algorithmic graph access is robust to input perturbation---a desirable
property for production deployment where schemas are diverse and
evolving.

\subsection{Model Differences}

Gemini and DeepSeek show consistent patterns but differ in magnitude.
DeepSeek-V3 benefits more from explicit graph context (GA outperforms
FT by 7.3pp for DeepSeek vs.\ only -1.1pp for Gemini), suggesting
that models with less domain exposure rely more heavily on structured
context injection.

\subsection{Positioning Against Stronger Agentic Methods}

Our baselines intentionally span the static-to-agentic spectrum rather
than exhaustively covering the agentic design space.  More sophisticated
approaches---planning-based agents with dynamic loop budgets (e.g.,
ToG-style multi-agent planners~\cite{sun2024tog}), hierarchical GraphRAG
with iterative subgraph expansion, or modular query decomposition
pipelines---may partially close the gap on compositional subtypes.
However, the convergence of TU (3~calls) and RC (5~code rounds) at
similar hard-question ceilings suggests that the bottleneck is the
LLM's compositional reasoning capacity rather than the tool interface.
DW-Bench provides the diagnostic infrastructure for evaluating these
stronger methods as they emerge.

\subsection{Limitations}

We do not enforce token-budget parity across baselines by design:
each paradigm receives context natural to its architecture (full schema
for FT, retrieved chunks for VR, BFS subgraph for GA, tool results
for TU).  45\% of questions exceed the GA 3-hop radius.  We evaluate
two frontier models; broader model coverage (GPT-4, open-weight LLMs)
would strengthen generality.  Tier~2 is limited to one synthetic dataset.
All lineage edges are assumed to form DAGs; cyclic ETL patterns
(e.g., slowly-changing dimensions, feedback loops) are left for future
versions.  Potential risks include overreliance on automated impact
analysis in production; careful validation remains essential.
